\begin{statementbox}

\textbf{\textcolor{CStatement}{条件}}
\begin{description}[style=nextline, leftmargin=2.8em, labelsep=0.5em, itemsep=0.35em]
    \item[\textbf{\textcolor{CStatement}{条件 1（非负变量）：}}] $x \geq 0$。
\end{description}

\textbf{\textcolor{CStatement}{结论}}
\begin{description}[style=nextline, leftmargin=2.8em, labelsep=0.5em, itemsep=0.45em]
    \item[\textbf{\textcolor{CStatement}{结论一（稳健双边）：}}]
    \textbf{\textcolor{CStatement}{直观描述：}} 对数函数被两个有理函数上下夹逼。
    \[
    \frac{2x}{x+2} \leq \ln(1+x) \leq \frac{x(x+2)}{2(x+1)}.
    \]

    \item[\textbf{\textcolor{CStatement}{结论二（标准下界）：}}]
    \textbf{\textcolor{CStatement}{直观描述：}} 对数函数大于等于一次项减二次修正项。
    \[
    \ln(1+x) \geq x - \frac{x^{2}}{2(1+x)}.
    \]

    \item[\textbf{\textcolor{CStatement}{推广结论（三阶上界）：}}]
    \textbf{\textcolor{CStatement}{直观描述：}} 更高精度的三次多项式上界。
    \[
    \ln(1+x) \leq x - \frac{x^{2}}{2} + \frac{x^{3}}{3}.
    \]
\end{description}

\textbf{\textcolor{CStatement}{等号/取等条件：}} 当且仅当 $x=0$ 时，所有不等式中的等号成立。

\end{statementbox}
